\rok{Фебруар 2024.}{B - Поправни тест другог теста}

Други поправни тест из Алгоритама и структура података

Први практични тест одржан 30.01.2024.

\zadatak{1}{Имплементација методе за изграђивање најкраће путање}
Написати имплементацију методе која треба да омогући креирање тренутно најкраће путање у оквиру Дајкстриног алгоритма.

\textbf{Додатне напомене за решавање задатка:}
\begin{itemize}
	\item У template-у је закоментарисана метода:
	\begin{Verbatim}[fontsize=\footnotesize, numbers=left, xleftmargin=10mm]
private static void BuildShortestPaths(int[] vertices, int[] previousVertices, HashSet<int>[] shortestPaths)
	\end{Verbatim}
	\item У Main методи обезбедити начин за тестирање и проверу нове функционалности.
\end{itemize}



\zadatak{2}{Дуплирање непонављајућих карактера}
Написати алгоритам временске комплексности $O(n)$ који ће из string-а који се прослеђује идентификовати карактере који се понављају, дуплирати све остале карактере улазног string-а, те тако измењени string исписати у конзоли.

\textbf{Додатне напомене за решавање задатка:}
\begin{itemize}
	\item Користити Hash табелу која је дата у шаблону.
	\item Имплементацију писати у оквиру закоментарисане методе:
	\begin{Verbatim}[fontsize=\footnotesize, numbers=left, xleftmargin=10mm]
public static string addMoreDuplicates(string str)
	\end{Verbatim}
	\item Алгоритам реализовати тако да постоји разлика између малих и великих карактера.
	\item За конверзију string-а у низ карактера користити методу \texttt{ToCharArray()} по принципу: \texttt{str.ToCharArray()}.
	\item Карактере није потребно експлицитно конвертовати у целобројни тип приликом убацивања у Hash табелу, али приликом каснијег исписивања користити експлицитно кастовање \texttt{(char)}.
	\item Редослед карактера приликом исписивања није битан.
\end{itemize}

\textbf{Примери:}
\begin{itemize}
	\item \textbf{Улаз 1:} \texttt{"Algoritam"}
	\item \textbf{Излаз 1:} \texttt{"AAllggoorriittaamm"} (ниједан карактер се не појављује више пута)
	\item \textbf{Улаз 2:} \texttt{"STRUKTURE podataka"}
	\item \textbf{Излаз 2:} \texttt{"SSTRUKKTUREE ppooddattakka"} (карактери \texttt{a}, \texttt{R}, \texttt{U} и \texttt{T} се појављују више пута)
\end{itemize}



\zadatak{3}{Провера путање са секвенцом чворова}
У оквиру стандардне структуре класе \texttt{BinarySearchTree} написати имплементацију методе која ће омогућити проверу постојања путање од корена до листа која се састоји из секвенце чворова одређене низом.

\textbf{Додатни захтеви за решавање задатка:}
\begin{itemize}
	\item Сматра се да путања са секвенцом чворова постоји уколико на путањи од корена до листа постоји секвенца чворова који имају вредности које одговарају вредностима у датом низу, и налазе се у редоследу који је одређен поретком у низу, при чему не може бити додатних чворова на путањи.
	\item У оквиру класе \texttt{BinarySearchTree} креирати јавну методу са потписом:
	\begin{Verbatim}[fontsize=\footnotesize, numbers=left, xleftmargin=10mm]
public bool ckeckIfPathExists(int[] arr)
	\end{Verbatim}
	\item Из претходно креиране јавне методе треба да се позива приватна метода са потписом:
	\begin{Verbatim}[fontsize=\footnotesize, numbers=left, xleftmargin=10mm]
private bool ckeckIfPathExists(Node curr, int[] arr, int index)
	\end{Verbatim}
	\item Приватна метода треба да садржи имплементацију алгоритма којим ће се омогућити провера секвенце вредности чворова на путањама од корена до листова.
	\item Алгоритам \textbf{ОБАВЕЗНО} писати у оквиру класе \texttt{BinarySearchTree}.
\end{itemize}

\textbf{Примери:}

\begin{center}
\begin{minipage}{\textwidth}
\centering
\begin{forest}
for tree={
	circle,
	draw,
	thick,
	fill=gray!20,
	minimum size=8mm,
	inner sep=1pt,
	s sep=8mm,
	l sep=12mm,
	edge={-, thick},
	font=\small
}
[8
	[4
		[2
			[, phantom]
			[3]
		]
		[5]
	]
	[9
		[, phantom]
		[11]
	]
]
\end{forest}

\vspace{0.3cm}
\textbf{Слика 1.} Пример BST-а.
\end{minipage}
\end{center}

\begin{itemize}
	\item \textbf{Улаз 1 за BST са слике 1:} \texttt{[8, 4, 2, 3]}
	\item \textbf{Излаз:} true
	\item \textbf{Улаз 2 за BST са слике 1:} \texttt{[8, 9, 10, 11]}
	\item \textbf{Излаз:} false
\end{itemize}

\begin{center}
\begin{minipage}{\textwidth}
\centering
\begin{forest}
for tree={
	circle,
	draw,
	thick,
	fill=gray!20,
	minimum size=8mm,
	inner sep=1pt,
	s sep=8mm,
	l sep=12mm,
	edge={-, thick},
	font=\small
}
[10
	[4
		[2
			[3]
			[, phantom]
		]
		[6
			[5]
			[9]
		]
	]
	[20
		[12
			[, phantom]
			[14]
		]
		[29
			[27]
			[, phantom]
		]
	]
]
\end{forest}

\vspace{0.3cm}
\textbf{Слика 2.} Пример BST-а.
\end{minipage}
\end{center}

\begin{itemize}
	\item \textbf{Улаз 1 за BST са слике 2:} \texttt{[10, 4, 6, 5]}
	\item \textbf{Излаз:} true
	\item \textbf{Улаз 2 за BST са слике 2:} \texttt{[10, 15, 14, 13]}
	\item \textbf{Излаз:} false
\end{itemize}



\sablon{Шаблон поправног теста за други тест фебруар 2024. група В}{2023/Februar/GrupaB/AISP2023_PopravniFebruar_Test2_GrupaB.linq}
